% Part: lambda-calculus
% Chapter: introduction
% Section: minimization

\documentclass[../../../include/open-logic-section]{subfiles}

\begin{document}

\olfileid{lam}{int}{min}
\olsection{$\lambd$-可定义函数在极小化运算下封闭}

\begin{lem}
设 $f(x,y)$ 是 $\lambd$-可定义的。令 $g$ 由
\[
g(x) \simeq \umin{y}{f(x,y)}.
\]
定义，则 $g$ 也是 $\lambd$-可定义的。
\end{lem}

\begin{proof}
思路大致如下。给定 $x$，我们将使用不动点 λ-项 $Y$ 来定义一个函数 $h_x(n)$，它从 $n$ 开始搜索 $y$；那么 $g(x)$ 就是 $h_x(0)$。函数 $h_x$ 可以表达为一个不动点方程的解：
\[
h_x(n) \simeq
\begin{cases}
n & \text{若 $f(x,n) = 0$} \\
h_x(n+1) & \text{否则。}
\end{cases}
\]

具体细节如下。由于 $f$ 是原始递归的，它由某个项 $F$ 所 $\lambd$-定义。回想一下，我们还有一个 λ-项 $D$，使得 $D(M, N, \bar{0}) \red M$ 且 $D(M, N, \bar{1})
\red N$。暂时固定 $x$。为了 $\lambd$-定义 $h_x$，我们需要找到一个依赖于 $x$ 的项 $H$，满足
\[
H(\num{n}) \equiv D(\num{n}, H(S(\num{n})), F(x, \num{n})).
\]
我们可以利用不动点项 $Y$ 做到这一点。首先，令 $U$ 为如下项
\[
\lambd[h][\lambd[z][D(z,(h(Sz)),F(x,z))]],
\]
再令 $H$ 为项 $YU$。注意，$H$ 中唯一的自由变元是 $x$。我们来证明 $H$ 满足上述方程。

根据 $Y$ 的定义，我们有
\[
H = YU \equiv U(YU) = U(H).
\]
特别地，对每个自然数 $n$，我们有
\begin{align*}
H(\num{n}) & \equiv U(H, \num{n}) \\
& \red D(\num{n}, H(S(\num{n})), F(x, \num{n})),
\end{align*}
如所求。请注意，如果在最后一行中将数码 $\num{m}$ 代入 $x$，那么当 $F(\num{m},
\num{n})$ 归约到 $\num{0}$ 时，该表达式归约到 $\num{n}$；而当 $F(\num{m}, \num{n})$ 归约到任何其他数码时，它归约到 $H(S(\num{n}))$。

证明的最后，令 $G$ 为 $\lambd[x][H(\num 0)]$。则 $G$ $\lambd$-定义了 $g$；换言之，对每个 $m$，若 $g(m)$ 有定义，则 $G(\num m)$ 归约到 $\overline {g(m)}$，否则 $G(\num m)$ 无范式。
\end{proof}

\end{document}